\documentclass[tikz,border=8pt]{standalone}
\usepackage{ctex}
\usepackage{amsmath}
\usepackage{pgfplots}
\pgfplotsset{compat=1.18}
\usetikzlibrary{calc,positioning,arrows.meta,decorations.markings}

%% 图 opt-p2-05-1：梯度下降轨迹
%% 椭圆等高线 + 迭代点 x0,x1,...x5 沿之字形路径收敛到中心
%% 标注步长 alpha 影响

\begin{document}
\begin{tikzpicture}[scale=1.1]

  % 坐标轴
  \draw[->, gray!60, thin] (-3.3, 0) -- (3.3, 0)
    node[right, font=\scriptsize] {$x_1$};
  \draw[->, gray!60, thin] (0, -2.5) -- (0, 2.7)
    node[above, font=\scriptsize] {$x_2$};

  % 椭圆等高线（f(x1,x2) = x1^2/4 + x2^2，条件数 kappa=4）
  % 等高线半轴 a=2r, b=r，r = 0.5, 0.9, 1.3, 1.7, 2.1
  \foreach \r/\lbl in {0.55/{$c_1$}, 1.0/{$c_2$}, 1.45/{$c_3$}, 1.9/{$c_4$}, 2.3/{$c_5$}}{
    \draw[gray!45, thin] (0,0) ellipse ({2*\r} and {\r});
  }
  % 等高线标签
  \node[font=\scriptsize, gray!70] at (0.8, 0.38) {$c_1$};
  \node[font=\scriptsize, gray!70] at (1.45, 0.68) {$c_2$};
  \node[font=\scriptsize, gray!70] at (2.1, 0.97) {$c_3$};

  % 极小值点（中心）
  \fill[red!80!black] (0,0) circle (3pt)
    node[below right, font=\small] {$\mathbf{x}^*$};

  % 梯度下降迭代轨迹（手动模拟 f(x)=x1^2/4+x2^2，步长 alpha=0.4）
  % 梯度：(x1/2, 2*x2)
  % 从 x0=(2.4, 1.2) 出发
  % x1 = x0 - 0.4*(1.2, 2.4) = (2.4-0.48, 1.2-0.96) = (1.92, 0.24)
  % x2 = x1 - 0.4*(0.96, 0.48) = (1.92-0.384, 0.24-0.192) = (1.536, 0.048)
  % 为了图示好看，用夸张的之字形轨迹
  \coordinate (X0) at (2.4, 1.2);
  \coordinate (X1) at (1.9, -0.85);
  \coordinate (X2) at (1.3, 0.62);
  \coordinate (X3) at (0.85, -0.42);
  \coordinate (X4) at (0.52, 0.28);
  \coordinate (X5) at (0.28, -0.15);

  % 轨迹线（红色实线箭头）
  \draw[->, thick, red!70!black, >=Stealth] (X0) -- (X1);
  \draw[->, thick, red!70!black, >=Stealth] (X1) -- (X2);
  \draw[->, thick, red!70!black, >=Stealth] (X2) -- (X3);
  \draw[->, thick, red!70!black, >=Stealth] (X3) -- (X4);
  \draw[->, thick, red!70!black, >=Stealth] (X4) -- (X5);
  \draw[->, thick, red!70!black, >=Stealth, dashed] (X5) -- (0.12, -0.06);

  % 迭代点标注
  \fill[blue!80!black] (X0) circle (2.5pt)
    node[above right, font=\small] {$\mathbf{x}_0$};
  \fill[blue!80!black] (X1) circle (2pt)
    node[right, font=\small] {$\mathbf{x}_1$};
  \fill[blue!80!black] (X2) circle (2pt)
    node[above right, font=\small] {$\mathbf{x}_2$};
  \fill[blue!80!black] (X3) circle (2pt)
    node[right, font=\scriptsize] {$\mathbf{x}_3$};
  \fill[blue!80!black] (X4) circle (2pt)
    node[above right, font=\scriptsize] {$\mathbf{x}_4$};
  \fill[blue!80!black] (X5) circle (2pt)
    node[right, font=\scriptsize] {$\mathbf{x}_5$};

  % 步长标注（在 x0->x1 的中点）
  \node[font=\scriptsize, red!60!black, fill=white, inner sep=1.5pt]
    at ($(X0)!0.5!(X1) + (-0.6, 0)$)
    {步长 $\alpha$};

  % 负梯度方向小箭头（在 x0 处）
  \draw[->, orange!80!black, thick, >=Stealth]
    (X0) -- ++(0.0, -0.55) node[right, font=\scriptsize] {$-\nabla f$};

  % 说明框
  \node[draw=gray!50, thin, fill=white, rounded corners=2pt,
        font=\scriptsize, inner sep=5pt, align=left]
    at (3.0, -1.9)
    {$\mathbf{x}_{k+1} = \mathbf{x}_k - \alpha\nabla f(\mathbf{x}_k)$\\
     之字形 $\Leftarrow$ 条件数 $\kappa>1$\\
     $\kappa = L/\mu = 4$（此处）};

  % 整体标题
  \node[font=\large\bfseries] at (0, 3.1)
    {梯度下降轨迹：椭圆等高线 + 之字形收敛};

\end{tikzpicture}
\end{document}
